\section{Discussion, limitations and conclusion}
\label{sec:discussion}

\paragraph{Recommendations for registries.} (1) Run the arms concurrently on the same variants and never compare against a stored baseline. (2) Register the stopping rule and use a test that is valid at every stopping time: the betting test with a futility margin in general, or the discordant-pair SPRT for pass/fail scores at $\delta=0$. (3) Treat admission as online multiple testing: give each exam an e-LOND level when it starts, and spend alpha over a skill's lineage so that resubmission does not pay. (4) Generate hidden variants from a committed seed; this controls leakage and also supplies the hundreds of paired units that certification needs. (5) Lead credentials with the anytime-valid lower bound, not the stopped estimate. (6) Publish both arms' transcripts and the registered parameters so the decision can be recomputed.

\paragraph{What the guarantee assumes.} Two kinds of assumption carry the guarantees, and they differ in strength. Within an exam, validity needs the conditional null of Eq.~\eqref{eq:H}, which holds under the design of Section~\ref{sec:setup}. These departures break it: state shared across variants (agent memory, response caching, batched calls that couple variants); a provider change during an exam that affects one arm more than the other; repeated seeds, which turn one variant's outcome into several; exclusion rules that depend on the arm or on the outcome; and any exam-level latent state that is not fixed before the exam, as in the coin example of Section~\ref{sec:setup}. Concurrent arms, fresh seeds, stateless calls and arm-blind exclusion rules registered in $\theta$ protect against each. E21 measures each departure for a useless skill (Table~\ref{tab:nullstress}, Appendix~\ref{app:more}). Mild versions leave the deployed grid gate near its target: a model change that affects both arms, 1 to 5\% repeated seeds or excluded skill-arm losses, and small shared effects give at most \NNsMildMax\%, against \NNsControl\% under the design. Strong versions do not: a model change that hits only the skill arm gives \NNsSkillArm\%, shared effects over blocks of ten variants \NNsBlock\%, 20\% repeats or exclusions \NNsRepeat\% and \NNsExclude\%, and an exam-level shift of the lift by $\pm0.05$ gives \NNsCoin\%. The fixed-$n$ test fails in the same cells, so these are failures of the design, not of sequential testing; each is prevented by the corresponding protocol rule. Across exams the conditions are much weaker: e-LOND and lineage spending need only that each exam's stopped wealth is an e-value given the level fixed at its start, and they hold under arbitrary dependence between exams.

\begin{table}[t]
\centering
\small
\caption{Threat model. \emph{Prevented}: the protocol makes the attack fail. \emph{Detectable}: a verifier can see it from the published record. \emph{Assumed}: the guarantee relies on an honest party. \emph{External}: needs enforcement outside the protocol.}
\label{tab:threats}
\begin{tabular}{p{0.36\linewidth}lp{0.42\linewidth}}
\toprule
Threat & Status & Mechanism or requirement \\
\midrule
Skill tailored to the hidden variants & Prevented & Skill digest logged before the commitment; seed mixes a later beacon \\
Proctor picks favourable variants & Prevented & Commitment precedes the beacon value \\
Stopping at a lucky moment & Prevented & Test valid at every stopping time (Proposition~\ref{prop:valid}) \\
Abandoning unfavourable exams & Prevented & Commitments stay visible, and the registry counts every committed exam, finished or not, in e-LOND's index and the lineage's attempt count; abandoning only removes admissions \\
Resubmitting until a pass & Prevented if linked & Lineage spending (Proposition~\ref{prop:lineage}) \\
Arm- or outcome-dependent exclusions & Detectable & Exclusion rule registered in $\theta$; excluded runs published \\
Skill copies the generator or its solver & Detectable in part & Coverage statistic; reviewing the skill text (Section~\ref{sec:res-leak}) \\
Choice of generator and margin & Assumed & Registry chooses them before the skill is seen; the credential certifies lift on that generator only \\
Proctor leaks the seed & Assumed & Threshold proctors or seed release inside trusted hardware \\
Transcripts not produced by the stated model & External & Attested execution; not provided by the protocol \\
Sybil authors evading lineage linking & External & Identity verification or stake \\
\bottomrule
\end{tabular}
\end{table}

Table~\ref{tab:threats} summarizes what the protocol protects. ``Verifiable'' in this paper means that a third party can recompute the decision from the published transcripts and commitments; it does not mean attested execution.

\paragraph{Limitations.} The entropy price of Theorem~\ref{thm:price} is the cost of the lineage guarantee alone; under the deployed level, in a registry past its first few dozen exams, online FDR control is the binding constraint and the price is small (E19). The validity results are theorems; the cost results depend on an outcome model whose nuisance parameters have no published values. We sweep them and fit them to the released SkillsBench moments, and the cost claims move within the ranges of Section~\ref{sec:res-cost}; the resampling estimates there assume that variants behave like the released configurations, which we cannot check, and they are optimistic because the skills were selected on observed lift. The reanalyses use benchmarks whose units (requirement batches, configurations) are not the hidden parametric variants we propose, and the SkillsBench registry lacks the self-generated condition. The generated-variant studies use two models; the benchmark study uses one. Nemotron's answers reached the token limit in \NLnCutPct\% of calls in studies 2 and 3 and \NLnFiveCutPct\% in study 5, which lowers the pass rates the gate sees; a reasoning model's budget is part of what a credential certifies. On the benchmark most tasks were excluded for harness failures, and studies 2 to 4 use task families that we designed. On the one family built by others (study 5) no published skill helped; the one text that did help, our harmful skill with Nemotron, was not certified within the cap. So we have shown the gate refusing on an external family but not admitting on one. When exams can be batched in windows and their data are independent, stopped e-BH \citep{wang2025stopped} is more powerful than e-LOND in our simulations (E20); our gate does not use it because a registry decides exams one at a time, and the lineage theory needs levels fixed at the start of each exam. The moment-matched outcome model's fitted correlation lies at the edge of the searched grid. None of the live studies tests whether any particular way of producing skills, such as peer synthesis, passes the gate more often. The verification algorithm establishes that the decision follows from the published transcripts, not that the transcripts were produced by the stated model; remote model APIs are not attestable end to end, and trusted execution has known side channels. A proctor who leaks the seed to an author defeats commit and reveal; threshold seeding by independent proctors is the natural fix and is not analysed here. Lineage alpha-spending is only as good as the registry's ability to link resubmissions, and the leakage experiment uses one synthetic task family.

\paragraph{Conclusion.} Skills are becoming the unit in which agent capability is shared, and registries already decide which ones to list. Those decisions are sequential, multiple and repeatable, and the per-skill evidence that benchmarks collect today is enough for a one-off fixed-sample analysis of some skills but certifies only a few once the analysis must survive sequential testing, a growing library and resubmission. A gate built from a paired exam on committed hidden variants, a betting test with a futility margin, e-LOND with lineage alpha-spending, and a recomputable dual-arm credential holds its error rates in every condition we tested that meets the design of Section~\ref{sec:setup}, and its cost can be planned: if hidden variants behave like the released configurations, from about fifty to a few hundred variants per certified skill.
